\chapter{Documented thesis case}
\label{app:tfg}

This appendix summarizes the assembly, software environment and results from the
thesis. They describe one configuration and should not be read as general AKD1000
requirements.

\section{Assembly used}

The thesis assembly consisted of the components listed in
\cref{tab:componentes-tfg} \cite{poveda_tfg_2026}.

\begin{longtable}{@{}L{0.27\textwidth}L{0.43\textwidth}L{0.22\textwidth}@{}}
\caption{Hardware used in the thesis.}
\label{tab:componentes-tfg}\\
\toprule
\textbf{Component} & \textbf{Identification} & \textbf{Notes} \\
\midrule
\endfirsthead

\toprule
\textbf{Component} & \textbf{Identification} & \textbf{Notes} \\
\midrule
\endhead

Host &
Raspberry Pi 5 Model B, 8~GB RAM, \code{Rev 1.1}. &
Purchased as part of a db-tronic kit. \\

Storage &
64~GB M.2 NVMe SSD. &
Included in the kit; the exact model was not retained. \\

Cooling and case &
Active cooling system and case included in the kit. &
The commercial part number was not retained. \\

Host power supply &
27~W power supply. &
Included in the kit. \\

Adapter &
Waveshare \emph{PCIe TO PCIe x1 Board (C)} for Raspberry Pi 5. &
Connected to the Raspberry Pi through FFC. \\

Cable &
FFC marked \code{PCIe-Cable-50mm}. &
50~mm long. \\

Board &
BrainChip Akida PCIe Board AKD1000. &
The device appeared as \code{BC.00.000.002}. \\

Board power &
12~V DC through the carrier. &
Powered separately from the Raspberry Pi. \\

\bottomrule
\end{longtable}

The physical chain was:

\begin{center}
Raspberry Pi 5
$\rightarrow$
FFC cable
$\rightarrow$
Waveshare carrier
$\rightarrow$
BrainChip AKD1000
\end{center}

The Raspberry Pi was powered through USB-C and the AKD1000 board through the carrier's
12~V input.

\section{Software environment}

\Cref{tab:tfg-env} lists the versions used during the final campaign.

\begin{table}[H]
\centering
\caption{Thesis software environment.}
\label{tab:tfg-env}
\begin{tabularx}{\textwidth}{@{}L{0.32\textwidth}Y@{}}
\toprule
\textbf{Component} & \textbf{Version} \\
\midrule
Operating system & Debian GNU/Linux 12 (bookworm) \\
Python & 3.11.2 \\
TensorFlow & 2.19.1 \\
Keras API & \code{tf\_keras} \\
Akida & 2.19.1 \\
CNN2SNN & 2.19.1 \\
QuantizeML & 1.2.3 \\
Device & \code{BC.00.000.002} \\
Host configuration & \code{performance} governor; \code{throttled=0x0} \\
\bottomrule
\end{tabularx}
\end{table}

These versions describe the environment used in the experiment. Fully reproducing
the results also requires the models, data and the same measurement procedure.

\section{AoA models}

The thesis used two networks to estimate the angle of arrival (\emph{Angle of Arrival},
AoA). Each sample contained eight integer envelope intensity values between 0 and 15.

The model output contained 180 values, one for each class. Class $c$ corresponded
to the angle

\begin{equation}
\theta = 2c.
\end{equation}

The two final models were:

\begin{table}[H]
\centering
\caption{AoA models used in the thesis.}
\label{tab:tfg-models}
\begin{tabularx}{\textwidth}{@{}L{0.23\textwidth}L{0.28\textwidth}L{0.23\textwidth}Y@{}}
\toprule
\textbf{Model} & \textbf{Architecture} & \textbf{Weights excluding biases} & \textbf{Use} \\
\midrule
Compact &
$8\rightarrow10\rightarrow180$ &
1,880 &
Low-complexity model \\

Main &
$8\rightarrow957\rightarrow180$ &
179,916 &
Higher-accuracy model \\
\bottomrule
\end{tabularx}
\end{table}

Angular error was calculated accounting for the circular nature of the angle:

\begin{equation}
e_{\mathrm{circ}}
=
((\hat{\theta}-\theta+180)\bmod 360)-180.
\end{equation}

For example, a prediction of $0^\circ$ for a reference of $358^\circ$ gives an error
of $2^\circ$, not $-358^\circ$.

The FBZ models and dataset used in the original campaign are not part of this release.
The test model included in the repository is used only to check the environment and board.

\section{Results}

Evaluation used 576 samples. The main model achieved an angular MAE of $4.09^\circ$,
compared with $9.13^\circ$ for the compact model \cite{poveda_tfg_2026}.

Among the configurations studied, the lowest estimated energy points were:

\begin{itemize}
  \item compact model: \code{Economy + HwPr}, with 0.565~ms,
        770.0~mW and 435.2~$\mu$J;
  \item main model: \code{Performance + AllNps}, with 0.756~ms,
        731.8~mW and 553.3~$\mu$J.
\end{itemize}

\begin{figure}[H]
  \centering
  \includegraphics[width=0.88\textwidth]{nine-mode-results.png}
  \caption[AoA model results]
  {Latency and estimated energy for both models across the nine clock and mapping
  mode combinations evaluated in the thesis.}
  \label{fig:tfg-results}
\end{figure}

Latency was measured from the host around the execution call. Power was measured
through the Akida API, rather than as total consumption of the Raspberry Pi, carrier
and board.

These results describe only the platform used in the thesis.
